\begin{frame}

\frametitle{Qestion3 (3)}

Answer whether or not the cnf-formula $\phi = (x_{35}\lor x_{36}\lor x_{37}) \land (\lnot x_{35}\lor\lnot x_{36}) \land (\lnot x_{36}\lor\lnot x_{37}) \land (\lnot x_{35}\lor\lnot x_{37})$ is satisfiable.  If the answer is yes, then also answer an example of the truth assignments that make $\phi=1$.  

左記の連言標準形ブール式 $\phi$ が充足可能か否か答えなさい．充足可能である場合はさらに，$\phi$の値を$1$にするような真理値割り当ての例を答えなさい．

\end{frame}

\begin{frame}

\frametitle{充足可能性}

ブール式$\phi$が充足可能とは，$\phi=1$となるような変数への$0$と$1$の割り当てが存在すること．

つまり，\textcolor{red}{$\phi=1$となる$x_{35}$，$x_{36}$，$x_{37}$の組み合わせを一つ以上挙げる}ことで証明可能．

\end{frame}

\begin{frame}

\frametitle{充足可能の証明}

連言標準形ブール式$\phi$は充足可能である．実際，

$x_{35}=1$，$x_{36}=0$，$x_{37}=0$のとき，

\begin{align}\displaystyle
\phi &= (1\lor0\lor0) \land (\lnot1\lor\lnot0) \land (\lnot0\lor\lnot0) \land (\lnot1\lor\lnot0) \notag \\
&= 1 \land 1 \land 1 \land 1 \notag \\
&= 1\notag
\end{align}

とである．

\end{frame}

\begin{frame}

\frametitle{論理値表の作成}

$\phi = (x_{35}\lor x_{36}\lor x_{37}) \land (\lnot x_{35}\lor\lnot x_{36}) \land (\lnot x_{36}\lor\lnot x_{37}) \land (\lnot x_{35}\lor\lnot x_{37})$

より，

$\phi=0 \Leftrightarrow (x_{35}\lor x_{36}\lor x_{37}=0) \text{または} (\lnot x_{35}\lor\lnot x_{36}=0) \text{または} (\lnot x_{36}\lor\lnot x_{37}=0) \text{または} (\lnot x_{35}\lor\lnot x_{37}=0)$

である．

ここで，$\phi=0$となる$x_{35}$，$x_{36}$，$x_{37}$の組み合わせを考える．

\end{frame}

\begin{frame}

\frametitle{論理値表の作成}

(i) $x_{35}\lor x_{36}\lor x_{37}=0$になる組み合わせは，

\[
\begin{array}{c|c|c|c}
x_{35} & x_{35} & x_{35} & \phi \\
\hline
0 & 0 & 0 & 0 \\
 &  &  &  \\
 &  &  &  \\
 &  &  &  \\
 &  &  &  \\
\end{array}
\]

\end{frame}

\begin{frame}

\frametitle{論理値表の作成}

(ii) $\lnot x_{35}\lor\lnot x_{36}=0$になる組み合わせは，

\[
\begin{array}{c|c|c|c}
x_{35} & x_{35} & x_{35} & \phi \\
\hline
0 & 0 & 0 & 0 \\
1 & 1 & 0 & 0 \\
 &  &  &  \\
 &  &  &  \\
1 & 1 & 1 & 0 \\
\end{array}
\]

\end{frame}

\begin{frame}

\frametitle{論理値表の作成}

(iii) $\lnot x_{36}\lor\lnot x_{37}=0$になる組み合わせは，

\[
\begin{array}{c|c|c|c}
x_{35} & x_{35} & x_{35} & \phi \\
\hline
0 & 0 & 0 & 0 \\
1 & 1 & 0 & 0 \\
0 & 1 & 1 & 0 \\
 &  &  &  \\
1 & 1 & 1 & 0 \\
\end{array}
\]

\end{frame}

\begin{frame}

\frametitle{論理値表の作成}

(iv) $\lnot x_{35}\lor\lnot x_{37}=0$になる組み合わせは，

\[
\begin{array}{c|c|c|c}
x_{35} & x_{35} & x_{35} & \phi \\
\hline
0 & 0 & 0 & 0 \\
1 & 1 & 0 & 0 \\
0 & 1 & 1 & 0 \\
1 & 0 & 1 & 0 \\
1 & 1 & 1 & 0 \\
\end{array}
\]

\end{frame}

\begin{frame}

\frametitle{論理値表の作成}

これで$\phi=0$となる全ての組み合わせが分かった．

これ以外の場合は，$\phi\neq0$，つまり$\phi=1$となる．したがって，

\[
\begin{array}{c|c|c|c}
x_{35} & x_{35} & x_{35} & \phi \\
\hline
1 & 0 & 0 & 1 \\
0 & 1 & 0 & 1 \\
0 & 0 & 1 & 1 \\
\end{array}
\]

が答えとなる．
\end{frame}
